\documentclass[a4paper,12pt]{article}
\usepackage{amssymb,amsmath,amsfonts,amsthm}
\usepackage{graphicx}
\usepackage{fullpage}
\usepackage{multirow}
\usepackage{colortbl}
\usepackage[latin1]{inputenc}
\usepackage[T1]{fontenc}
\usepackage[brazil]{babel}
\usepackage{calligra}
\usepackage{accents}
\usepackage{cancel}

% Definindo o estilo de numera��o dos subteoremas
\newtheorem{theorem}{Teorema}[section]

% Definindo um novo estilo espec�fico para subteoremas com a numera��o desejada
\newtheorem{subtheorem}{Teorema}[theorem]

% Customizando o contador do subtheorem
\renewcommand{\thesubtheorem}{\thesection.\arabic{theorem}(\alph{subtheorem})}

\begin{document}
	
% Define a cor da linha de cancelamento como vermelha
\renewcommand{\CancelColor}{\color{red}}

\setcounter{section}{14}

\setcounter{theorem}{1}

\setcounter{subtheorem}{0}

\begin{subtheorem}
	Se $X\sim \text{Geom�trica}(p)$, a m�dia  � dada por $ E(X) = \frac{1}{p} $.
\end{subtheorem}

\begin{proof}
	Temos que
	$$
		\begin{aligned}
			E(X) &= \sum\limits_{x=1}^\infty x P(X=x)\\
			&= \sum\limits_{x=1}^\infty x (1-p)^{x-1}p\\
			&= \sum\limits_{x=1}^\infty (x-1+1) (1-p)^{x-1}p\\
			&= \sum\limits_{x=1}^\infty \left[ (x-1) (1-p)^{x-1}p + (1-p)^{x-1}p\right]\\
			&= \sum\limits_{x=1}^\infty (x-1) (1-p)^{x-1}p + \sum\limits_{x=1}^\infty (1-p)^{x-1}p\\
			&= \sum\limits_{x=1}^\infty (x-1) (1-p)^{x-1}p + 1\\
		\end{aligned}
	$$
	Fazendo $y=x-1$, temos
	$$
		\begin{aligned}
			E(X) &= \sum\limits_{y=0}^{\infty} y(1-p)^yp + 1\\
			&= (1-p)\sum\limits_{y=1}^{\infty} y(1-p)^{y-1}p + 1\\
			&= (1-p)E(X) + 1\\
		\end{aligned}
	$$
	Ou seja,
	$$
		\begin{aligned}
			E(X) &= (1-p)E(X) + 1\\
			E(X) - (1-p)E(X) &= 1\\
			\cancel{E(X)} - \cancel{E(X)} + pE(X) &= 1\\
			E(X) &= \frac{1}{p}
		\end{aligned}
	$$
\end{proof}

\begin{subtheorem}
	Se $X\sim \text{Geom�trica}(p)$, a vari�ncia  � dada por $ Var(X) = \frac{1-p}{p^2} $.
\end{subtheorem}
	
\begin{proof}
	Sabemos que a vari�ncia � definida como $Var(X) = E(X^2) - [E(X)]^2$. Vamos calcular primeiro $E(X^2)$.
	$$
	\begin{aligned}
		E(X^2) &= \sum\limits_{x=1}^\infty x^2 P(X=x)\\
		&= \sum\limits_{x=1}^\infty x^2 (1-p)^{x-1} p\\
		&= \sum\limits_{x=1}^\infty (x -1 +1)^2 (1-p)^{x-1} p\\
		&= \sum\limits_{x=1}^\infty [(x -1)^2 + 2(x-1) +1] (1-p)^{x-1} p\\
		&= \sum\limits_{x=1}^\infty (x -1)^2 (1-p)^{x-1} p + \sum\limits_{x=1}^\infty 2(x-1) (1-p)^{x-1} p + \sum\limits_{x=1}^\infty (1-p)^{x-1} p
	\end{aligned}
	$$
	Fazendo $y=x-1$ e como $\sum\limits_{x=1}^\infty (1-p)^{x-1} p = 1$, temos
	$$
	\begin{aligned}
		E(X^2) &= \sum\limits_{y=0}^\infty y^2 (1-p)^y p + \sum\limits_{y=0}^\infty 2y (1-p)^y p + 1\\
		&= 0 + \sum\limits_{y=1}^\infty y^2 (1-p)^y p + 0 + \sum\limits_{y=1}^\infty 2y (1-p)^y p + 1\\
		&= (1-p)\sum\limits_{y=1}^\infty y^2 (1-p)^{y-1} p + 2(1-p)\sum\limits_{y=1}^\infty y (1-p)^{y-1} p + 1\\
		&= (1-p)E(X^2) + 2(1-p)E(X) + 1\\
	\end{aligned}
	$$
	Usando $E(X)=\frac{1}{p}$, temos
	$$
	\begin{aligned}
		E(X^2) &= (1-p)E(X^2) + 2(1-p)\frac{1}{p} +1\\
		E(X^2) - (1-p)E(X^2) &= \frac{2(1-p) + p}{p}\\
		\cancel{E(X^2)} - \cancel{E(X^2)} + pE(X^2) &= \frac{2 -2p + p}{p}\\
		E(X^2) &= \frac{2 -p}{p^2}\\
	\end{aligned}
	$$
	o que d� o resultado
	$$
		Var(X) = E(X^2) - [E(X)]^2 = \frac{2 -p}{p^2} - \left[\frac{1}{p}\right]^2 = \frac{2 -p -1}{p^2} = \frac{1-p}{p^2}
	$$
\end{proof}	
	
%%%%%%%%%%%%%%%%%%%%%%%%%%%%%%%%%%%%%%%%%%%%%%%%%%%%%%%%%%%%%%%%%%%%%%
\newpage

\setcounter{theorem}{2}

\setcounter{subtheorem}{0}

\begin{subtheorem}
	Se $X\sim \text{Binomial Negativa}(r, p)$, a m�dia  � dada por $ E(X) = \frac{r}{p} $.
\end{subtheorem}

\begin{proof}
	Temos que
	$$
	\begin{aligned}
		E(X) &= \sum\limits_{x=r}^\infty x P(X=x)\\
		&= \sum\limits_{x=r}^\infty x \binom{x-1}{r-1}p^r(1-p)^{x-r}\\
		&= \sum\limits_{x=r}^\infty x \frac{(x-1)!}{(x-r)!(r-1)!}p^r(1-p)^{x-r}\\
		&= \sum\limits_{x=r}^\infty \frac{x!}{(x-r)!(r-1)!}\left[\frac{r}{r}\right]p^r(1-p)^{x-r}\\
		&= \sum\limits_{x=r}^\infty \frac{x!}{(x-r)!\,r!}\,r \,p^r(1-p)^{x-r}\\
		&= \sum\limits_{x=r}^\infty \binom{x}{r} r  \left[\frac{p}{p}\right] p^r(1-p)^{x-r}\\
		&= \sum\limits_{x=r}^\infty \binom{x}{r} \frac{r}{p} \, p^{r+1}(1-p)^{x-r}\\
		&= \frac{r}{p} \sum\limits_{x=r}^\infty \binom{x}{r}  p^{r+1}(1-p)^{x-r}\\
	\end{aligned}
	$$
	Fazendo $x=y-1$, temos
	$$
	\begin{aligned}
		E(X) &= \frac{r}{p} \sum\limits_{y=r+1}^\infty \binom{y-1}{r}  p^{r+1}(1-p)^{(y-1)-r}
	\end{aligned}
	$$
	Agora fazendo $r=r+1$, temos
	$$
	\begin{aligned}
		E(X) &= \frac{r}{p} \sum\limits_{y=s}^\infty \binom{y-1}{s-1}  p^{s}(1-p)^{(y-1)-(s-1)}\\
		&= \frac{r}{p} \sum\limits_{y=s}^\infty \binom{y-1}{s-1}  p^{s}(1-p)^{y-s}\\
		&= \frac{r}{p} \sum\limits_{y=s}^\infty P(Y=y), \text{ com } Y\sim \text{Binomial Negativa}(s, p)\\
		&= \frac{r}{p}\cdot 1\\
		&= \frac{r}{p}
	\end{aligned}
	$$
\end{proof}

\begin{subtheorem}
	Se $X\sim \text{Binomial Negativa}(r, p)$, a vari�ncia  � dada por $ Var(X) = \frac{r(1-p)}{p^2} $.
\end{subtheorem}

\begin{proof}
	Sabemos que a vari�ncia � definida como $Var(X) = E(X^2) - [E(X)]^2$. Vamos calcular primeiro $E(X^2)$.
	$$
	\begin{aligned}
		E(X^2) &= \sum\limits_{x=1}^\infty x^2 P(X=x)\\
		&= \sum\limits_{x=r}^\infty x^2 \binom{x-1}{r-1}p^r(1-p)^{x-r}\\
		&= \sum\limits_{x=r}^\infty x^2 \frac{(x-1)!}{(x-r)!(r-1)!} p^r(1-p)^{x-r}\\
		&= \sum\limits_{x=r}^\infty x^2 \frac{(x-1)!}{(x-r)!(r-1)!}\left[\frac{rp}{rp}\right] p^r(1-p)^{x-r}\\
		&= \sum\limits_{x=r}^\infty x \frac{x(x-1)!}{(x-r)!\,r(r-1)!}\left[\frac{r}{p}\right] pp^r(1-p)^{x-r}\\
		&= \frac{r}{p} \sum\limits_{x=r}^\infty x \frac{x!}{(x-r)!\,r!} p^{r+1}(1-p)^{x-r}\\
		&= \frac{r}{p} \sum\limits_{x=r}^\infty x \binom{x}{r} p^{r+1}(1-p)^{x-r}\\
	\end{aligned}
	$$
	Fazendo $x=y-1$, temos
	$$
	\begin{aligned}
		E(X^2) &= \frac{r}{p} \sum\limits_{y=r+1}^\infty (y-1) \binom{y-1}{r} p^{r+1}(1-p)^{(y-1)-r}
	\end{aligned}
	$$
	Fazendo agora $s=r+1$, obtemos
	$$
	\begin{aligned}
		E(X^2) &= \frac{r}{p} \sum\limits_{y=s}^\infty (y-1) \binom{y-1}{s-1} p^{s}(1-p)^{(y-1)-(s-1)}\\
		&= \frac{r}{p} \sum\limits_{y=s}^\infty (y-1) \binom{y-1}{s-1} p^{s}(1-p)^{y-s}\\
		&= \frac{r}{p} \left[ \sum\limits_{y=s}^\infty y \binom{y-1}{s-1} p^{s}(1-p)^{y-s} - \sum\limits_{y=s}^\infty \binom{y-1}{s-1} p^{s}(1-p)^{y-s}\right]\\
		&= \frac{r}{p} \left[E(Y) - \sum\limits_{y=s}^\infty P(Y=y) \right],
	\end{aligned}
	$$
	em que $Y\sim \text{Binomial Negativa}(s, p)$. Logo, $E(Y) = \frac{s}{p} = \frac{r+1}{p}$ e $\sum\limits_{y=s}^\infty P(Y=y) =1$. Da�,
	$$
		E(X^2) = \frac{r}{p}\left[\frac{r+1}{p} - 1\right] = \frac{r}{p}\left[\frac{r+1-p}{p}\right] = \frac{r^2 + r -rp}{p^2}
	$$
	Portanto,
	$$
	\begin{aligned}
		Var(X) &= E(X^2) - [E(X)]^2\\
		&= \frac{r^2 + r -rp}{p^2} - \left[\frac{r}{p}\right]^2\\
		&= \frac{\cancel{r^2}+r-rp -\cancel{r^2}}{p^2}\\
		&=\frac{r(1-p)}{p^2}
	\end{aligned}
	$$
\end{proof}


\end{document}